% =====================================================================
% RQ3.  Every number here is produced by make_router_table.py and
% matches tab:router exactly.  Do not edit numbers by hand.
% =====================================================================

\paragraph{RQ3: Is the decision to intervene learnable from what a single call reveals?}
\label{sec:rq3-router}

Section~\ref{sec:rq2} showed that a static intervention helps or hurts depending on whether the
item needed external knowledge in the first place, and that applying one indiscriminately is
significantly worse than not applying it at all. That turns the problem from a retrieval-quality
question into a decision question: can a system tell, before paying for a second call, whether
that call will repair its answer or break it.

\paragraph{Protocol.} The gate is a logistic model over $27$ features fitted on the $5{,}110$
public items. Its regularisation strength and decision threshold are chosen by five-fold
cross-validation inside that split, the fitted coefficients and threshold are written to disk, and
only then is the held-out feature matrix constructed. A whitelist guard raises if a held-out
identifier reaches the fitting path, so the separation is enforced rather than asserted. The
held-out half is scored once. No feature reads the gold answer, any correctness flag, or any
statistic computed over held-out items.

\paragraph{The gate sees one call, not three.} This is the design choice that makes the comparison
meaningful. A selector whose features include the predictions of every condition has already paid for
every condition, so it sits in the same budget class as majority voting over those conditions, and majority
voting requires no training. We therefore restrict the gate to what the Direct call itself
exposes, its sandbox execution and parser status, its error code, its token counts and latency,
and the option letter it chose, together with properties of the item text. The gate pays for a
second call only where it fires, which is $1.08\times$ the compute of Always Direct on the
held-out half and $1.04\times$ on its context-complete part.

\paragraph{Result.} Table~\ref{tab:router} reports every policy at its true budget. All three
fixed policies that always intervene are significantly worse than never intervening, Function-RAG
at $-1.16$ points ($p=0.007$) and \textsc{FunctionGraph-RAG} at $-1.10$ ($p=0.020$). The gate is
the only policy that significantly improves on Always Direct, gaining $0.39$ points on the
held-out half ($p=0.045$, bootstrap $95\%$ CI $[+0.04,+0.77]$) and $0.55$ points on the
context-complete part ($p<0.001$, CI $[+0.26,+0.83]$). On the context-complete part it fires on
$4.3\%$ of items and converts $31$ items while breaking $8$, a rescue to harm ratio of almost four
to one. Neither trivial baseline reproduces this. Switching whenever the two conditions disagree is
identical to always intervening, and switching at random at the gate's own firing rate gains
$0.16$ points on the held-out half ($p=0.302$) and nothing at all on the context-complete part,
which shows the gain comes from the decision and not from the rate.

\paragraph{What this does and does not establish.} The effect is small in absolute terms and we do
not claim otherwise. Its significance rests on a favourable rescue to harm ratio on a small number
of interventions rather than on a large accuracy movement. What it establishes is that the
intervention decision carries signal observable at inference time, which the aggregate results of
Section~\ref{sec:rq2} would not have predicted, and that this signal is available from a single
call rather than requiring every branch to be executed first. The three-condition post-hoc selector,
which does require that, reaches $71.36$ but is not separable from Always Direct ($p=0.107$) and
is barely above majority voting at the same budget, $71.36$ against $71.21$.

\paragraph{The ceiling this leaves.} The oracle over the same three conditions reaches $77.75$, and that
number is not a stronger method but an identity: it equals $100$ minus the $22.25\%$ of held-out
items on which all three conditions fail. The routing decision is live on only $784$ of $5{,}088$ items,
since on the remaining $84.6\%$ every condition agrees and the choice cannot matter, and within that
space $432$ items where Direct succeeds and an intervention would break it outnumber the $352$
where the reverse holds. A gate must therefore clear $44.9\%$ precision merely to break even. The
$6.92$ points separating Always Direct from the oracle bound what any selector over these conditions can
recover. The larger quantity is elsewhere: the $1{,}132$ items no condition answers correctly are worth
$22.25$ points, of which $440$ lack the context needed to answer at all and $692$ are fully
specified and simply not solved. Those are not selection problems, and closing them requires
enlarging the action space rather than choosing better within it.
